% page 48 du fac-simile (image p-047 du scan)
% bloc 160.0 x 242.0 mm ; marge gauche 8.0 mm
\pgc{-12.700mm}{0.000mm}{
\l{\cel{8}40}
\l{}
\l{\sou{arkeologo}.\cel{2}La persono qua su okupas pri arkeologio.}
\l{}
\l{\sou{arkesporo}.\cel{1}\sou{(bot.)}\cel{2}Parto interna di sporuyo che la muski,}
\l{\cel{3}filiki. - DEFIS.}
\l{}
\l{\sou{arki-}\cel{3}Prefixo qua indikas \textquotedbl{}en grado eminenta\textquotedbl{} ; lu, do,}
\l{\cel{3}indikas grado supera. - DEFIRS.}
\l{}
\l{\sou{arkimandrito}. (\sou{Eklezio Greka}). Superioro di monakeyo unesma-}
\l{\cel{3}klasa. - DEFIRS.}
\l{}
\l{\sou{arkipelago.}\cel{2}Grupo de insuli. - DEFIRS.}
\l{}
\l{\sou{arkiplasmo}. (\sou{biol}.) Substanco ye qua konsistas la direkto-}
\l{\cel{3}sferi en kariokinezo. - DEFIS.}
\l{}
\l{\sou{arkitekto}. La persono qua facas la mapo di edifico e direktas}
\l{\cel{3}la konstrukto di olca. (\sou{anke metaf.})\cel{2}- DEFIRS.}
\l{}
\l{\sou{arkitekturo}.\cel{2}La arto konstruktar edifici. - DEFIRS.}
\l{}
\l{\sou{arkitravo}.\cel{2}Infra parto di la entablamento qua su apogas sur}
\l{\cel{3}la kapiteli di la koloni. - DEFIRS.}
\l{}
\l{\sou{arkivo}.\cel{2}Singla kompozanto ek la ensemblo de la akti, doku-}
\l{\cel{3}menti pri la historio di populo, di urbo, di instituto pu-}
\l{\cel{3}blika o privata, di komerco-societo. - DEFIRS.}
\l{}
\l{\sou{arkivolto}.\cel{2}Kadro di la arkado di pordo, di fenestro. -DEFIRS.}
\l{}
\l{\sou{arkonto}.\cel{2}Un ek la non oficieri elektita singlayare, qui gu-}
\l{\cel{3}vernis la republiko di la urbo Atena,antique. - DEFIRS.}
\l{}
\l{\sou{arkoplasmo.}-\sou{(biol.}) Atrakto-sfero o direkto-vezikulo, judikata}
\l{\cel{3}kom un ek la organi prima di la celulo. - DEFIS.}
\l{}
\l{\sou{arktika}.\cel{2}Situita en la mi-sfero di la Terglobo, qua regardas}
\l{\cel{3}vers la stelaro \textquotedbl{}Granda Urso\textquotedbl{}, vers la nordo-polo di la}
\l{\cel{3}Terglobo.- DEFIRS.}
\l{}
\l{\sou{arlekino}. Persono komika di la farso Italiana ; lu aparas sur}
\l{\cel{3}la ceneyo, kun vesto konsistanta ye peci triangula, diversa-}
\l{\cel{3}kolora, e kun ligna sabro en la manuo. - DEFIRS.}
\l{}
\l{\sou{armo}.\cel{2}Instrumento uzata da la homo por atakar o por su}
\l{\cel{3}defensar. - DEFIS.}
\l{}
\l{\sou{armadilo}. (\sou{zool.}) Mamifero sen-denta, kelke simila a la aselo.}
\l{\cel{3}- L.\cel{1}\sou{dasypus}. - DEFIS.}
\l{}
\l{\sou{armatoro}. Posedanto ed equipanto di navo. - FIRS.}
\l{}
\l{\sou{armaturo}. Asembluro de peci de fero, por mantenar kunligite}
\l{\cel{3}la diversa parti di karpentajo o di masonajo. - DEFIRS.}
}
